\documentclass{article}

\usepackage[T1]{fontenc}

\usepackage{booktabs,graphicx,siunitx}

\begin{document}

\section{Kruskal-Wallis}

\begin{tabular}{lllll}
\toprule
Group & N & Median & IQR & Mean rank \\
\midrule
control & 20 & 33.55 & 9.63 & 24.90 \\
drug\_a & 20 & 36.50 & 11.97 & 28.15 \\
drug\_b & 20 & 38.70 & 6.60 & 38.45 \\
\bottomrule
\end{tabular}

\begin{tabular}{llll}
\toprule
H & df & p & $\varepsilon$\textsuperscript{2} [95\% CI] \\
\midrule
6.56 & 2 & .038 & 0.11 [0.03, 1.00] \\
\bottomrule
\end{tabular}

\begin{tabular}{lll}
\toprule
Pair & Z & padj \\
\midrule
control - drug\_a & 0.59 & .556 \\
control - drug\_b & 2.45 & .042 \\
drug\_a - drug\_b & 1.87 & .124 \\
\bottomrule
\end{tabular}

\includegraphics{figures/01_kruskal-wallis/figure_boxplot.png}

The nonparametric counterpart to one-way ANOVA. The H/p tests whether any group tends to have systematically higher or lower values (stochastic dominance) --- read it as a "median difference" only when the groups have similar distribution shapes. If significant, Dunn's post-hoc shows which pairs differ, with adjusted p-values. $\varepsilon$\textsuperscript{2} is the effect size. Your significance threshold ($\alpha$) is 0.05.

A Kruskal-Wallis test gave H(2)=6.56, p = .038, $\varepsilon$\textsuperscript{2}=.11 [.03, 1.00].

Statistical basis: Kruskal-Wallis H test (nonparametric counterpart to one-way ANOVA) (Kruskal, W. H., Wallis, W. A. (1952). "Use of ranks in one-criterion variance analysis." Journal of the American Statistical Association, 47(260), 583-621.); Dunn's post-hoc pairwise comparisons (Kassambara A (2023). rstatix: Pipe-Friendly Framework for Basic Statistical Tests. R package.); $\varepsilon$\textsuperscript{2} effect size (Ben-Shachar MS, Lüdecke D, Makowski D (2020). "effectsize: Estimation of Effect Size Indices and Standardized Parameters." Journal of Open Source Software, 5(56), 2815.).

\end{document}